% Evaluación Comparativa de Ejecución de Código vs Llamadas Directas
% en Sistemas Agénticos Basados en MCP
% LNCS format - Springer
%
\documentclass[runningheads]{llncs}
%
\usepackage[T1]{fontenc}
\usepackage[spanish]{babel}
\usepackage{graphicx}
\usepackage{booktabs}
\usepackage{amsmath}
\usepackage{multirow}
\usepackage{xcolor}
\usepackage{listings}
\usepackage{hyperref}
\renewcommand\UrlFont{\color{blue}\rmfamily}
\urlstyle{rm}
%
% Listings config for Python code
\lstset{
  language=Python,
  basicstyle=\ttfamily\small,
  keywordstyle=\color{blue},
  commentstyle=\color{gray},
  stringstyle=\color{red},
  breaklines=true,
  frame=single,
  numbers=left,
  numberstyle=\tiny\color{gray}
}
%
\begin{document}
%
\title{Evaluaci\'on Comparativa de Ejecuci\'on de C\'odigo vs Llamadas Directas en Sistemas Ag\'enticos Basados en MCP}
%
\titlerunning{Ejecuci\'on de C\'odigo vs Llamadas Directas en MCP}
%
% TODO: Fill in author information
\author{Autor Uno\inst{1}}
%
\authorrunning{A. Uno}
%
% TODO: Fill in institution
\institute{Universidad, Ciudad, Pa\'is\\
\email{autor@universidad.edu}}
%
\maketitle
%
% ============================================================================
% RESUMEN
% ============================================================================
\begin{abstract}
% TODO: Escribir después de completar las secciones técnicas (~250 palabras)
% Estructura: contexto → problema → método → resultados clave → conclusión
Los sistemas agénticos basados en modelos de lenguaje grande (LLMs) utilizan
el Model Context Protocol (MCP) para acceder a herramientas externas. Existen
dos paradigmas de ejecución: llamadas directas a herramientas MCP y generación
de código ejecutado en sandbox. Este trabajo presenta una evaluación comparativa
sistemática de ambos enfoques mediante un diseño factorial 2$\times$2 que combina
arquitectura (agente único vs multi-agente) con modo de ejecución (MCP directo
vs ejecución de código), evaluado sobre tres tipos de tareas de complejidad
creciente utilizando ocho servidores MCP. Se miden tokens consumidos, latencia,
precisión y tasa de éxito. Los resultados muestran\ldots
% TODO: Completar con resultados reales

\keywords{MCP \and ejecuci\'on de c\'odigo \and tool calling \and sistemas multi-agente \and eficiencia de tokens}
\end{abstract}
%
% ============================================================================
% 1. INTRODUCCIÓN (Write 3rd)
% Maps to thesis: Chapter 1 (Contexto, Planteamiento, Motivación, Objetivos)
% ============================================================================
\section{Introducci\'on}
\label{sec:intro}

% 1.1 Contexto: LLMs, agentes, MCP como estándar emergente
% TODO: Párrafo sobre auge de LLMs y sistemas agénticos
% TODO: Párrafo sobre MCP como protocolo estándar para tool use

% 1.2 Problema: dos paradigmas de ejecución, trade-offs desconocidos
% TODO: Párrafo describiendo MCP directo vs code execution
% TODO: Párrafo sobre la falta de evaluaciones comparativas

% 1.3 Motivación: costos, latencia, necesidad de guías
% TODO: Párrafo sobre implicaciones prácticas (costos de tokens, latencia)

% 1.4 Contribuciones
% TODO: Lista de contribuciones:
% - Framework experimental (OrchAIstra)
% - Evaluación sistemática 2×2 × 3 task types
% - Análisis cuantitativo de métricas
% - Guías para selección de modo de ejecución

% 1.5 Estructura del artículo
% TODO: Párrafo describiendo organización de secciones

%
% ============================================================================
% 2. ANTECEDENTES Y TRABAJO RELACIONADO (Write 2nd)
% Maps to thesis: Chapter 2 (Revisión de Literatura) + Chapter 3 (Marco Teórico)
% ============================================================================
\section{Antecedentes y Trabajo Relacionado}
\label{sec:background}

\subsection{Model Context Protocol (MCP)}
\label{sec:mcp}
% TODO: Origen y especificación del protocolo MCP (Anthropic, Nov 2024)
% TODO: Arquitectura cliente-servidor, transporte stdio/SSE
% TODO: Ciclo de vida: discovery → invocation → response
% TODO: Adopción: Claude Desktop, AWS Strands, editores de código

\subsection{Ejecuci\'on de C\'odigo en Sistemas Ag\'enticos}
\label{sec:code-exec}
% TODO: Patrón de code execution (Anthropic engineering blog)
% TODO: Sandboxing: Docker, E2B, Cloudflare Workers
% TODO: Comparación con tool calling directo
% TODO: Strands code_interpreter como referencia

\subsection{Arquitecturas Multi-Agente}
\label{sec:multi-agent}
% TODO: Topologías: orquestador, enjambre, pipeline
% TODO: Frameworks: AutoGen, CrewAI, Strands
% TODO: Delegación de tareas y coordinación

\subsection{Trabajo Relacionado}
\label{sec:related}
% TODO: Benchmarks de tool-use existentes
% TODO: Evaluaciones de eficiencia de tokens
% TODO: Comparaciones de modos de ejecución (si existen)
% TODO: Diferenciación de nuestra contribución

%
% ============================================================================
% 3. DISEÑO EXPERIMENTAL (Write 1st - most concrete)
% Maps to thesis: Chapter 4 (Diseño Experimental) + Chapter 5 (Metodología)
% ============================================================================
\section{Dise\~no Experimental}
\label{sec:design}

\subsection{Matriz Factorial}
\label{sec:matrix}
% TODO: Tabla 2×2 (arquitectura × modo de ejecución)
% A1: Agente único + MCP directo
% A2: Agente único + Ejecución de código
% A3: Multi-agente + MCP directo
% A4: Multi-agente + Ejecución de código

% Placeholder for the factorial table
% \begin{table}[h]
% \caption{Matriz factorial 2$\times$2 de configuraciones experimentales.}\label{tab:matrix}
% \centering
% \begin{tabular}{lcc}
% \toprule
%  & MCP Directo & Ejecución de Código \\
% \midrule
% Agente Único & A1 & A2 \\
% Multi-Agente & A3 & A4 \\
% \bottomrule
% \end{tabular}
% \end{table}

\subsection{Suite de Tareas}
\label{sec:tasks}
% TODO: Tres tipos de tareas con complejidad creciente
% T1: Tarea simple (1 herramienta) - morse, cocktails
% T2: Tarea encadenada (pipeline secuencial) - nutrition, alphavantage, nasa
% T3: Tarea compleja (30+ herramientas) - devtoolkit, lastfm

% TODO: Tabla con tipos de tarea, descripción, servidores representativos

\subsection{Servidores MCP}
\label{sec:servers}
% TODO: Descripción de los 8 servidores
% TODO: Tabla: nombre, #herramientas, tipo de tarea, API externa
% TODO: Explicar implementación dual (MCP server + Python tools)

\subsection{Framework: OrchAIstra}
\label{sec:framework}
% TODO: Arquitectura del framework de evaluación
% TODO: Diagrama de componentes (AgentFactory, DockerSandbox, MetricsCollector)
% TODO: Flujo de ejecución de un experimento
% TODO: Sistema de evaluación por LLM (grading)

\subsection{M\'etricas}
\label{sec:metrics}
% TODO: Definición formal de cada métrica
% - Tokens: input, output, total
% - Latencia: total_latency_ms
% - Precisión: LLM-graded score (0-100)
% - Tool calls: count, success_rate
% - Failure modes: categorización de errores

\subsection{Hip\'otesis}
\label{sec:hypotheses}
% TODO: Hipótesis formales
% H1: Code execution reduce tokens para T2 y T3
% H2: MCP directo tiene menor latencia para T1
% H3: Precisión equivalente entre modos
% H4: Existe punto de cruce en número de herramientas

\subsection{Variables de Control}
\label{sec:controls}
% TODO: Modelo LLM fijo, temperatura, número de repeticiones
% TODO: Configuración de Docker (memory, CPU, network)

%
% ============================================================================
% 4. RESULTADOS (Write after experiments)
% Maps to thesis: Chapter 6 (Resultados)
% ============================================================================
\section{Resultados}
\label{sec:results}

\subsection{Resultados por Arquitectura}
\label{sec:results-arch}
% TODO: Tablas comparativas A1-A4 para cada tipo de tarea
% TODO: Métricas agregadas por configuración

\subsection{An\'alisis de Tokens}
\label{sec:results-tokens}
% TODO: Gráficas de consumo de tokens por configuración y tarea
% TODO: Desglose input vs output tokens
% TODO: Comparación MCP directo vs code execution

\subsection{An\'alisis de Latencia}
\label{sec:results-latency}
% TODO: Tiempos de ejecución por configuración
% TODO: Overhead de Docker en code execution
% TODO: Latencia per-tool vs end-to-end

\subsection{An\'alisis de Precisi\'on}
\label{sec:results-precision}
% TODO: Tasas de éxito por configuración y tarea
% TODO: Distribución de scores del LLM grader
% TODO: Análisis de modos de fallo

\subsection{An\'alisis Estad\'istico}
\label{sec:results-stats}
% TODO: Tests de significancia (t-test, ANOVA)
% TODO: Intervalos de confianza
% TODO: Tamaño del efecto

%
% ============================================================================
% 5. DISCUSIÓN (Write after results)
% Maps to thesis: Chapter 7 (Modelo de Decisión) + Chapter 8 (Discusión)
% ============================================================================
\section{Discusi\'on}
\label{sec:discussion}

\subsection{Interpretaci\'on de Resultados}
\label{sec:interpretation}
% TODO: Relación con hipótesis H1-H4
% TODO: Patrones observados

\subsection{An\'alisis de Modos de Fallo}
\label{sec:failures}
% TODO: Cuándo falla cada enfoque y por qué
% TODO: Categorización de errores

\subsection{Implicaciones Pr\'acticas}
\label{sec:implications}
% TODO: Guías para desarrolladores
% TODO: Cuándo usar MCP directo vs code execution
% TODO: Consideraciones de costo y rendimiento

\subsection{Limitaciones}
\label{sec:limitations}
% TODO: Limitaciones del estudio
% TODO: Amenazas a la validez

%
% ============================================================================
% 6. CONCLUSIONES (Write last)
% Maps to thesis: Chapter 9 (Conclusiones)
% ============================================================================
\section{Conclusiones}
\label{sec:conclusions}
% TODO: Resumen de hallazgos principales
% TODO: Contribuciones del trabajo
% TODO: Trabajo futuro: más modelos, más tareas, optimización de prompts

%
% ============================================================================
% ACKNOWLEDGMENTS & DISCLOSURES
% ============================================================================
\begin{credits}
\subsubsection{\ackname}
% TODO: Agradecimientos

\subsubsection{\discintname}
Los autores declaran no tener conflictos de inter\'es relevantes al contenido
de este art\'iculo.
\end{credits}
%
% ============================================================================
% REFERENCES
% ============================================================================
%
\bibliographystyle{splncs04}
% \bibliography{references}
%
% Inline bibliography until .bib file is ready
\begin{thebibliography}{20}

% MCP Protocol
\bibitem{mcp_spec}
Anthropic: Model Context Protocol Specification.
\url{https://modelcontextprotocol.io/} (2024)

% MCP Python SDK
\bibitem{mcp_sdk}
Anthropic: MCP Python SDK.
\url{https://github.com/modelcontextprotocol/python-sdk} (2024)

% Anthropic code execution
\bibitem{anthropic_code_exec}
Anthropic: Code execution with MCP. Anthropic Engineering Blog (2025)

% AWS Strands
\bibitem{strands}
AWS: Strands Agents -- Code Interpreter Tool.
\url{https://github.com/strands-agents/tools} (2025)

% Docker security
\bibitem{docker_security}
Docker Inc.: Docker Engine Security.
\url{https://docs.docker.com/engine/security/} (2024)

% E2B
\bibitem{e2b}
E2B: Cloud Runtime for AI Agents.
\url{https://github.com/e2b-dev/e2b} (2024)

% TODO: Add more references as sections are written
% - LLM agent frameworks (AutoGen, CrewAI)
% - Tool-use benchmarks
% - Token efficiency studies
% - Multi-agent architectures

\end{thebibliography}
%
\end{document}
